\chapter{Conceptual Framework and
Methodology}\label{conceptual-framework-and-methodology}

This chapter bridges the theoretical framework established in Chapter 3
and the formal analysis to be presented in the next chapters. It
presents the conceptual framework that operationalises the theoretical
concepts for IBC specifically, followed by a step-by-step methodology
for conducting the formal specification, LTS extraction, and behavioural
comparison.

The chapter addresses the following questions:

\begin{itemize}
    \item How are the theoretical concepts (atomic actions, operators, bisimulation) instantiated for IBC?
    \item What is the step-by-step procedure for producing the process algebra specification?
    \item How are LTS models extracted from process algebra terms?
    \item How is the unified/stacked/conflicting classification determined?
\end{itemize}

\section{Conceptual Framework}\label{conceptual-framework}

The theoretical framework in Chapter 3 introduced general concepts from
process algebra and LTS theory. This subsection defines those concepts
specifically for the IBC protocol.

\subsection{Observable Behaviour in
IBC}\label{observable-behaviour-in-ibc}

For the purposes of this research, \textbf{observable behaviour} of IBC
consists of actions that are visible to an external observer of the
protocol. The following actions are considered observable:

\begin{itemize}
    \item \textbf{Packet send events:} When a source chain commits a packet to a channel
    \item \textbf{Packet receive events:} When a target chain accepts a packet
    \item \textbf{Acknowledgement events:} When a packet is successfully processed
    \item \textbf{Timeout events:} When a packet fails to be delivered within the deadline
    \item \textbf{Handshake events:} OpenInit, OpenTry, OpenAck, OpenConfirm
    \item \textbf{Channel closure events:} CloseInit, CloseConfirm
\end{itemize}

The following are considered \textbf{internal (silent) actions} and are
abstracted as \(\tau\):

\begin{itemize}
    \item Relayer packet forwarding between chains
    \item Cryptographic proof generation and verification (Merkle proofs)
    \item Light client state synchronisation
    \item Internal state updates within a single blockchain
\end{itemize}

\subsection{Trace Definition for IBC}\label{trace-definition-for-ibc}

A \textbf{trace} in IBC is a finite sequence of observable actions from
the initial channel establishment to either successful termination or
deadlock. For example, a successful data transfer trace is:

\[
Trace_{data} = \langle OpenInit, OpenTry, OpenAck, OpenConfirm, Send, Receive, Ack \rangle
\]

A trace with timeout is:

\[
Trace_{timeout} = \langle OpenInit, OpenTry, OpenAck, OpenConfirm, Send, Timeout \rangle
\]

\subsection{Branching Structure in
IBC}\label{branching-structure-in-ibc}

Branching structure occurs when the protocol has a choice between two or
more continuations from a given state. In IBC, branching occurs at:

\begin{itemize}
    \item \textbf{Acknowledgement vs Timeout:} After sending a packet, the protocol may either receive an acknowledgement (\(ack\)) or expire (\(timeout\)). This is modelled as \(ack + timeout\).
    \item \textbf{Handshake acceptance vs rejection:} During channel establishment, the target chain may accept (\(OpenTry\)) or reject (\(OpenTryReject\)).
\end{itemize}

\subsection{Deadlock in IBC}\label{deadlock-in-ibc}

A \textbf{deadlock} in IBC is a state from which no further progress is
possible but which is not a successful termination. Examples include:

\begin{itemize}
    \item \textbf{Incomplete handshake:} OpenInit sent but no OpenTry received, and timeout mechanism not triggered
    \item \textbf{Trapped packet:} Packet sent but neither acknowledged nor timed out due to relayer failure
    \item \textbf{Channel identifier exhaustion:} No available channel identifiers for new connections
\end{itemize}

These deadlock states are captured in the LTS as states with no outgoing
transitions that are not marked as termination states (\(\checkmark\)).

\section{Research Paradigm: Formal Modelling as
Inquiry}\label{research-paradigm-formal-modelling-as-inquiry}

This research adopts a \textbf{formal modelling paradigm} rather than an
empirical one. The following distinctions are critical:

\begin{table}[h]
\centering
\begin{tabular}{|l|l|}
\hline
\textbf{Empirical Research} & \textbf{This Research (Formal Modelling)} \\
\hline
Observes real system execution & Constructs mathematical abstraction \\
Measures performance or behaviour & Derives behaviour from specification \\
Statistical validity & Correctness of mapping validity \\
Requires live network or implementation & Works from protocol standards (ICS documents) \\
Generalises from samples & Generalises from formal proof \\
\hline
\end{tabular}
\caption{Comparison of Empirical and Formal Modelling Paradigms}
\end{table}

Validity in this research is established through:

\begin{enumerate}
    \item \textbf{Face validity:} The process algebra specification accurately reflects the IBC standards (ICS-004, ICS-020) as written
    \item \textbf{Systematic mapping:} Every IBC action has a corresponding atomic action in the alphabet
    \item \textbf{Complete SOS application:} All LTS transitions are derived using the SOS rules from Chapter 3
    \item \textbf{Reproducibility:} The complete specification is provided in the appendix for independent verification
\end{enumerate}

\subsection{The Three Analytical Scenarios (Unified, Stacked,
Conflicting)}\label{the-three-analytical-scenarios-unified-stacked-conflicting}

The central question of this research is whether IBC\textquotesingle s
asset-transfer and data-transfer behaviours are Unified, Stacked, or
Conflicting. Table 4.1 provides the operational definitions for each
scenario using the theoretical framework from Chapter 3.

\begin{table}[h]
\centering
\begin{tabular}{|p{3cm}|p{5cm}|p{5cm}|}
\hline
\textbf{Scenario} & \textbf{Operational Definition} & \textbf{LTS Property to Verify} \\
\hline
\textbf{Unified} & Asset and data behaviours are weakly bisimilar: \(P_{asset} \approx P_{data}\) & LTS of asset and data have identical branching structure; every observable action in one has a matching action in the other after \(\tau\) abstraction \\
\hline
\textbf{Stacked} & Asset behaviour contains all data traces but adds constraints: \(Traces(P_{data}) \subseteq Traces(P_{asset})\) but \(P_{asset} \not\approx P_{data}\) & All data traces appear as subsequences of asset traces; asset traces include additional lock/mint actions; branching structures differ \\
\hline
\textbf{Conflicting} & Parallel composition introduces deadlocks not present in isolation: \(Deadlock(P_{asset} \parallel P_{data}) \not\subseteq Deadlock(P_{asset}) \cup Deadlock(P_{data})\) & Running asset and data transfers simultaneously over the same channels creates deadlock states that do not occur when only asset or only data transfers are executed \\
\hline
\end{tabular}
\caption{Operational Definitions for Unified, Stacked, and Conflicting}
\end{table}

\subsection{Conceptual Framework
Diagram}\label{conceptual-framework-diagram}

Figure 4.1 illustrates the overall conceptual framework, showing the
flow from IBC standards to the final classification.

\begin{figure}[htbp]
\centering
\resizebox{0.95\textwidth}{!}{%
\begin{tikzpicture}[
    node distance=0.6cm and 0.8cm,
    box/.style={rectangle, draw, text width=3cm, align=center, 
                font=\scriptsize, inner sep=3pt, rounded corners=2pt},
    arrow/.style={->, thick, >=stealth},
    stage/.style={font=\bfseries\small, anchor=south}
]

% Stage labels
\node[stage] at (0.5, 1.2) {Phase 1};
\node[stage] at (4.2, 1.2) {Phase 2};
\node[stage] at (7.9, 1.2) {Phase 3};
\node[stage] at (11.6, 1.2) {Phase 4};
\node[stage] at (15.3, 1.2) {Phase 5};

% Row 1 (main flow)
\node[box] (A) {IBC Standards\\ICS-004, ICS-020};
\node[box, right=of A] (B) {Domain Analysis\\Atomic Actions \& $\gamma$};
\node[box, right=of B] (C) {Process Algebra\\$P_{ICS004}, P_{ICS020}$};
\node[box, right=of C] (D) {LTS Extraction\\States \& Transitions};
\node[box, right=of D] (E) {Behavioural\\Analysis};
\node[box, right=of E] (F) {Weak Bisimulation\\$P_{004} \approx P_{020}$?};
\node[box, right=of F] (G) {Classification\\Unified/Stacked/Conflicting};

% Arrows between main boxes
\draw[arrow] (A) -- (B);
\draw[arrow] (B) -- (C);
\draw[arrow] (C) -- (D);
\draw[arrow] (D) -- (E);
\draw[arrow] (E) -- (F);
\draw[arrow] (F) -- (G);

% Row 2 (detailed outputs below each phase)
\node[box, below=0.4cm of B, text width=2.5cm, font=\tiny] (Bdet) {Alphabet $Act$\\Ports $P$, Data $D$};
\node[box, below=0.4cm of C, text width=2.5cm, font=\tiny] (Cdet) {$P_{ICS004}$\\$P_{ICS020}$\\$P_{Combined}$};
\node[box, below=0.4cm of D, text width=2.5cm, font=\tiny] (Ddet) {LTS$_D$, LTS$_A$\\LTS$_C$};
\node[box, below=0.4cm of E, text width=2.5cm, font=\tiny] (Edet) {Traces\\Branching\\Deadlocks};
\node[box, below=0.4cm of F, text width=2.5cm, font=\tiny] (Fdet) {Trace sets\\Cones \& Foci\\$D_{combined}$};
\node[box, below=0.4cm of G, text width=2.5cm, font=\tiny] (Gdet) {Final\\Result};

\draw[arrow, dashed] (B) -- (Bdet);
\draw[arrow, dashed] (C) -- (Cdet);
\draw[arrow, dashed] (D) -- (Ddet);
\draw[arrow, dashed] (E) -- (Edet);
\draw[arrow, dashed] (F) -- (Fdet);
\draw[arrow, dashed] (G) -- (Gdet);

\end{tikzpicture}%
}
\caption{Conceptual Framework for IBC Behavioural Analysis (Five Phases)}
\label{fig:conceptual-framework}
\end{figure}

\section{Research Methodology}\label{research-methodology}

This research will employ a
\textbf{formal specification and analytical verification} design. The
methodology follows a top-down decomposition approach:

\begin{enumerate}
    \item Start with the IBC protocol standards (ICS-004, ICS-020) as the authoritative source of truth
    \item Decompose each standard into its constituent atomic actions and interaction patterns
    \item Translate these patterns into process algebra terms using the operators defined in Chapter 3
    \item Extract LTS models from the process algebra terms using SOS rules
    \item Analyse the LTS models to identify traces, branching structures, deadlocks, and synchronisation patterns
    \item Compare the asset and data behaviours using weak bisimulation and deadlock detection
    \item Classify the relationship as Unified, Stacked, or Conflicting
\end{enumerate}

\textbf{Tools to be used:}

\begin{itemize}
    \item Manual process algebra specification (ACP/CCS style notation)
    \item Manual LTS extraction and analysis (guided by SOS rules)
    \item Optional: Validation using process algebra toolset (mCRL2 or CADP) if available
\end{itemize}

\subsection{Phase 1: Domain Analysis and Atomic Action
Identification}\label{phase-1-domain-analysis-and-atomic-action-identification}

This phase addresses \textbf{Objective 1} and provides the foundation
for \textbf{RQ1}.

\begin{enumerate}
\item
  Decompose IBC Standards

  Obtain the authoritative IBC standards:

  \begin{itemize}
      \item ICS-004: Channel and Packet Semantics (data transport layer)
      \item ICS-020: Fungible Token Transfer (asset transfer layer)
  \end{itemize}

  Sources: Cosmos IBC documentation and specification repositories
  \citep{cosmos2024ibcintro} \citep{cosmos_ibc}.
\item
  Identify Communication Ports

  Define the set of communication ports \(P\) for IBC:

  \begin{itemize}
      \item \(p_{chainA}\): Port on source blockchain
      \item \(p_{chainB}\): Port on target blockchain
      \item \(p_{relayer}\): Port for off-chain relayer (internal, will be abstracted)
      \item \(p_{channel}\): Logical channel endpoint
  \end{itemize}
\item
  Define Data Types

  Define the set of data \(D\) for IBC packets:

  \begin{itemize}
      \item \(d_{packet}\): General packet payload (bytes, payload-agnostic)
      \item \(d_{asset}\): Asset transfer payload (amount, denomination, sender, receiver)
      \item \(d_{ack}\): Acknowledgement data (success/failure, optional return data)
      \item \(d_{channel\_id}\): Channel identifier
      \item \(d_{sequence}\): Packet sequence number
      \item \(d_{timeout}\): Timeout height and timestamp
  \end{itemize}
\item
  Define Atomic Actions

  Define the atomic actions for IBC:

  \textbf{Handshake actions (ICS-004):} \[
  \begin{aligned}
  & chan\_open\_init_{A}(id) \\
  & chan\_open\_try_{B}(id) \\
  & chan\_open\_ack_{A}(id) \\
  & chan\_open\_confirm_{B}(id)
  \end{aligned}
  \]

  \textbf{Packet actions (ICS-004):} \[
  \begin{aligned}
  & send\_packet_{A}(p,d) \\
  & recv\_packet_{B}(p,d) \\
  & ack_{B}(p,d_{ack}) \\
  & recv\_ack_{A}(p,d_{ack}) \\
  & timeout_{A}(p)
  \end{aligned}
  \]

  \textbf{Asset-specific actions (ICS-020):} \[
  \begin{aligned}
  & lock_{A}(asset) \\
  & burn_{A}(asset) \quad \text{(alternative to lock)} \\
  & mint_{B}(voucher) \\
  & verify\_proof_{B}(proof)
  \end{aligned}
  \]
\item
  Define Communication Function

  Define the communication function \(\gamma\) that pairs send and read
  actions:

  \[
  \begin{aligned}
  & \gamma(send\_packet_{A}(p,d), recv\_packet_{B}(p,d)) = c\_packet(p,d) \\
  & \gamma(ack_{B}(p,d_{ack}), recv\_ack_{A}(p,d_{ack})) = c\_ack(p,d_{ack}) \\
  & \gamma(chan\_open\_init_{A}(id), chan\_open\_try_{B}(id)) = c\_open(id) \\
  & \gamma(chan\_open\_ack_{A}(id), chan\_open\_confirm_{B}(id)) = c\_confirm(id)
  \end{aligned}
  \]

  Mismatched or isolated actions (send without matching read, or vice
  versa) are blocked by the encapsulation operator \(\partial_H\).
\item
  Deliverable

  A complete alphabet specification containing:

  \begin{itemize}
      \item Set of ports \(P\)
      \item Set of data types \(D\)
      \item Set of atomic actions \(Act = Act_{handshake} \cup Act_{packet} \cup Act_{asset}\)
      \item Communication function \(\gamma\)
      \item Set of internal actions \(I\) to be abstracted (relayer forwarding, proof verification)
  \end{itemize}
\end{enumerate}

\subsection{Phase 2: Process Algebra
Specification}\label{phase-2-process-algebra-specification}

This phase addresses \textbf{Objective 2} and completes \textbf{RQ1}.

\begin{enumerate}
\item
  Specify ICS-004 (Data Transport)

  Define the channel handshake process: \[
  Handshake = chan\_open\_init_{A} \cdot (chan\_open\_try_{B} \cdot chan\_open\_ack_{A} \cdot chan\_open\_confirm_{B} + TimeoutReject)
  \]

  Define the packet flow process: \[
  PacketFlow = send\_packet_{A} \cdot (recv\_packet_{B} \cdot ack_{B} \cdot recv\_ack_{A} + timeout_{A})
  \]

  Define the full channel process (recursive for multiple packets): \[
  Channel = Handshake \cdot (PacketFlow)^* \cdot Close
  \] where \((PacketFlow)^*\) denotes zero or more repetitions of
  PacketFlow, and \(Close\) terminates the channel.

  Define the complete IBC data transport system: \[
  IBC_{Data} = \partial_H( Chain_A \parallel Relayer \parallel Chain_B )
  \] where:

  \begin{itemize}
      \item \(Chain_A\) and \(Chain_B\) are blockchain processes running Channel
      \item \(Relayer\) forwards packets internally (modelled as \(\tau\) actions)
      \item \(\partial_H\) blocks unmatched send/read actions (requires handshake completion)
  \end{itemize}
\item
  Specify ICS-020 (Asset Transfer)

  Define the lock/burn process on source chain: \[
  LockProcess = lock_{A}(asset) \cdot verify\_lock_{A} \cdot packet\_prep
  \]

  Define the mint/claim process on target chain: \[
  MintProcess = verify\_proof_{B} \cdot mint_{B}(voucher) \cdot ack\_send
  \]

  Define the full asset transfer (wrapped in ICS-004 packet): \[
  AssetTransfer = LockProcess \cdot SendPacket \cdot MintProcess + Rollback
  \] where \(Rollback\) handles timeout or verification failure by
  unlocking the asset.

  Define the complete IBC asset transfer system: \[
  IBC_{Asset} = \partial_H( ChainA_{asset} \parallel Relayer \parallel ChainB_{asset} )
  \]
\item
  Specify Combined System (for Conflict Detection)

  To test for conflicts, define the parallel composition of asset and
  data flows over the same channels: \[
  IBC_{Combined} = \partial_H( (Chain_A \parallel ChainA_{asset}) \parallel Relayer \parallel (Chain_B \parallel ChainB_{asset}) )
  \]
\item
  Deliverable

  Complete process algebra specifications for:

  \begin{itemize}
      \item \(P_{ICS004}\) (data transport only)
      \item \(P_{ICS020}\) (asset transfer only)
      \item \(P_{Combined}\) (both running concurrently)
  \end{itemize}

  All specifications follow the syntax defined in Chapter 3 and will be
  included in the appendices.
\end{enumerate}

\subsection{Phase 3: LTS Extraction}\label{phase-3-lts-extraction}

This phase addresses \textbf{Objective 3} and \textbf{RQ2}.

\begin{enumerate}
\item
  Generate LTS from Process Terms

  Using the SOS rules defined in Chapter 3 (Section 3.6), generate the
  LTS for each process term:

  \textbf{Procedure:}

  \begin{enumerate}
      \item Start from the initial state of the process term
      \item For each state, examine all possible SOS rule applications
      \item For each applicable rule, derive the next state and transition label
      \item Add the new state to the set \(S\) and the transition to \(\rightarrow\)
      \item Repeat until no new states are reachable (fixpoint)
      \item Mark termination states as \(\checkmark\)
      \item Mark deadlock states (no outgoing transitions, not terminated) as \(\delta\)
  \end{enumerate}
\item
  Record LTS Properties

  For each LTS (\(LTS_{Data}\), \(LTS_{Asset}\), \(LTS_{Combined}\)),
  record:

  \begin{enumerate}
      \item \textbf{States (S):} List of all reachable states with unique identifiers
      \item \textbf{Transitions (→):} Labelled edges between states
      \item \textbf{Partial traces:} All finite sequences of observable actions
      \item \textbf{Completed traces:} Sequences ending in \(\checkmark\) or \(\delta\)
      \item \textbf{Branching points:} States with multiple outgoing transitions (choice points)
      \item \textbf{Deadlock states:} States with no outgoing transitions and not marked \(\checkmark\)
      \item \textbf{Synchronisation patterns:} Transitions labelled with communication actions \(c_p(d)\)
  \end{enumerate}
\item
  Visualise LTS

  Produce transition graphs for each LTS. Example structure for
  \(LTS_{Data}\):

  \begin{verbatim}
  State S0 (Initial) --chan_open_init_A--> S1
  State S1 --chan_open_try_B--> S2
  State S2 --chan_open_ack_A--> S3
  State S3 --chan_open_confirm_B--> S4 (Channel Open)
  State S4 --send_packet_A--> S5
  State S5 --recv_packet_B--> S6
  State S6 (Choice) --ack_B--> S7 or --timeout_A--> S8
  ...
  \end{verbatim}
\item
  Deliverable

  For each specification (\(P_{ICS004}\), \(P_{ICS020}\),
  \(P_{Combined}\)):

  \begin{itemize}
      \item Complete LTS with enumerated states and transitions
      \item Set of partial traces
      \item Set of completed traces
      \item Set of deadlock states (if any)
      \item Branching structure documentation
  \end{itemize}
\end{enumerate}

\subsection{Phase 4: Behavioural
Comparison}\label{phase-4-behavioural-comparison}

This phase addresses \textbf{Objectives 4 and 5} and answers
\textbf{RQ3}.

\begin{enumerate}
\item
  Compare Trace Sets

  Compute and compare trace sets: \[
  \begin{aligned}
  T_{data} &= Traces(P_{ICS004}) \\
  T_{asset} &= Traces(P_{ICS020}) \\
  \end{aligned}
  \]

  Check:

  \begin{enumerate}
      \item \(T_{data} \subseteq T_{asset}\)? (Are all data traces also asset traces?)
      \item \(T_{asset} \subseteq T_{data}\)? (Are all asset traces also data traces?)
      \item \(T_{data} = T_{asset}\)? (Are the trace sets identical?)
  \end{enumerate}

  \textbf{Interpretation:}

  \begin{itemize}
      \item If \(T_{data} = T_{asset}\) → possible Unified (but must check branching)
      \item If \(T_{data} \subset T_{asset}\) → possible Stacked
      \item If neither inclusion holds → likely Conflicting or behaviours fundamentally different
  \end{itemize}
\item
  Check Weak Bisimulation

  Apply the Cones and Foci methodology \cite{Fokkink2004} to determine
  if \(P_{ICS004} \approx P_{ICS020}\):

  \textbf{Procedure:}

  \begin{enumerate}
      \item Construct a state mapping function \(\phi\) from states of \(P_{ICS020}\) to states of \(P_{ICS004}\)
      \item For each state in \(P_{ICS020}\), identify the matching "focus" state in \(P_{ICS004}\)
      \item Prove that for each visible action, the branching structure is preserved
      \item Prove that \(\tau\) steps in one process can be matched by zero or more \(\tau\) steps in the other
  \end{enumerate}

  If \(\phi\) exists and all proof obligations are satisfied, then
  \(P_{ICS004} \approx P_{ICS020}\) (weak bisimilar).
\item
  Detect Conflicts

  Compute deadlock sets: \[
  \begin{aligned}
  D_{data} &= Deadlock(P_{ICS004}) \\
  D_{asset} &= Deadlock(P_{ICS020}) \\
  D_{combined} &= Deadlock(P_{ICS004} \parallel P_{ICS020})
  \end{aligned}
  \]

  Check if: \[
  D_{combined} \not\subseteq D_{data} \cup D_{asset}
  \]

  If YES, then new deadlocks appear only when asset and data behaviours
  run concurrently. For each new deadlock, document the trace that leads
  to it.
\item
  Assign Final Classification

  Apply the decision rules from Table 4.1:
\item
  Deliverable

  \begin{itemize}
      \item Trace inclusion relationships (with evidence)
      \item Weak bisimulation proof or disproof (with Cones and Foci mapping)
      \item Deadlock sets for all three processes
      \item Final classification: \textbf{Unified}, \textbf{Stacked}, or \textbf{Conflicting}
      \item If Conflicting: Documented traces leading to new deadlocks
  \end{itemize}
\end{enumerate}

\subsection{Phase 5: Documentation and
Reusability}\label{phase-5-documentation-and-reusability}

This phase addresses \textbf{Objective 6}.

\begin{enumerate}
\item
  Standardise Notation

  Legends for all process algebra symbols used:

  \begin{table}[h]
  \centering
  \begin{tabular}{|l|l|}
  \hline
  \textbf{Symbol} & \textbf{Meaning} \\
  \hline
  \(a\) & Atomic action \\
  \(\cdot\) & Sequential composition \\
  \(+\) & Alternative composition (choice) \\
  \(\parallel\) & Parallel composition \\
  \(\partial_H\) & Encapsulation (block actions in set \(H\)) \\
  \(\tau_I\) & Abstraction (rename actions in \(I\) to \(\tau\)) \\
  \(\tau\) & Silent (internal) action \\
  \(\delta\) & Deadlock \\
  \(\checkmark\) & Successful termination \\
  \(\gamma\) & Communication function \\
  \(s_p(d)\) & Send action on port \(p\) with data \(d\) \\
  \(r_p(d)\) & Read action on port \(p\) with data \(d\) \\
  \(c_p(d)\) & Communication action (synchronised send/receive) \\
  \(\approx\) & Weak bisimulation \\
  \(\leftrightarrow_{rb}\) & Branching bisimulation \\
  \hline
  \end{tabular}
  \caption{Process Algebra Notation Legend}
  \end{table}
\item
  Document Assumptions and Abstractions

  List of all abstractions made (from Scope and Limitations in Chapter
  1):

  \begin{enumerate}
      \item Cryptographic verification (Merkle proofs, light client sync) is abstracted as \(\tau\)
      \item Economic incentives (gas costs, slashing, liquidity constraints) are not modelled
      \item Relayer behaviour is abstracted as \(\tau\) (only observable effects are packet delivery)
      \item Only ICS-004 and ICS-020 are modelled; ICS-27 (Interchain Accounts) and ICS-31 (Interchain Queries) are excluded
      \item Security properties beyond observable behaviour (e.g., double-spending resistance) are not verified
      \item No empirical validation or live protocol execution is performed
  \end{enumerate}
\item
  Produce Reusable Specification

  Format the complete specification as an appendix with:

  \begin{itemize}
      \item Complete alphabet definition
      \item Complete process algebra terms for \(P_{ICS004}\), \(P_{ICS020}\), and \(P_{Combined}\)
      \item Communication function \(\gamma\) definition
      \item Encapsulation set \(H\) and abstraction set \(I\)
      \item Commentary explaining each part of the specification
  \end{itemize}

  This specification is designed to be reusable for:

  \begin{itemize}
      \item Future comparison with Polkadot XCM
      \item Future comparison with Axelar General Message Passing (GMP)
      \item Input to automated process algebra tools (mCRL2, CADP, or similar)
  \end{itemize}
\item
  Deliverable

  \begin{itemize}
      \item An appendix for complete and commented process algebra specification
      \item An appendix for the extracted LTS models (states and transitions)
      \item Documentation of assumptions for future researchers
  \end{itemize}
\end{enumerate}

\section{Validation and Quality
Criteria}\label{validation-and-quality-criteria}

Unlike empirical research, formal modelling research requires different
validation criteria.

\subsection{Correctness of
Formalisation}\label{correctness-of-formalisation}

\begin{itemize}
    \item \textbf{Face validity:} The process algebra specification is reviewed against ICS-004 and ICS-020 documents line by line. Every protocol action in the standards has a corresponding atomic action or process term in the specification.
    \item \textbf{Systematic mapping:} A mapping table documents how each IBC standard element translates to process algebra.
    \item \textbf{Expert review (if available):} A researcher with formal methods or blockchain interoperability expertise reviews the specification for accuracy.
\end{itemize}

\subsection{Completeness of LTS
Extraction}\label{completeness-of-lts-extraction}

\begin{itemize}
    \item \textbf{State space coverage:} The LTS includes all states reachable from the initial state via SOS rule application. No reachable state is omitted.
    \item \textbf{Termination handling:} All termination states (\(\checkmark\)) are correctly identified.
    \item \textbf{Deadlock detection:} All deadlock states are identified and documented.
\end{itemize}

\subsection{Reproducibility}\label{reproducibility}

\begin{itemize}
    \item The complete process algebra specification is provided in the appendices.
    \item Each step of LTS generation is documented (transition-by-transition)
    \item If tools are used (mCRL2, CADP), the input scripts and commands are included
    \item The methodology is described with sufficient detail for another researcher to replicate
\end{itemize}

\subsection{Limitations}\label{limitations}

\begin{itemize}
    \item \textbf{Abstraction may hide subtle bugs:} Cryptographic failures (e.g., Merkle proof forgery) are not modelled; economic attacks (e.g., relayer bribery) are not captured; timing-dependent issues may be lost in \(\tau\) abstraction
    \item \textbf{Manual LTS extraction:} Manual derivation may introduce human error (mitigated by systematic SOS application and cross-checking)
    \item \textbf{Scalability:} State explosion may occur for highly concurrent scenarios (mitigated by compositional analysis and abstraction)
    \item \textbf{No empirical validation:} Results are theoretical predictions, not guarantees about live IBC implementations
\end{itemize}

\section{Ethical Considerations}\label{ethical-considerations}

This research involves no human subjects, no personal data collection,
and no live blockchain interaction. The formal modelling and analysis
are conducted entirely on publicly available protocol specifications
(ICS-004, ICS-020). No economic or security risks are introduced as no
actual assets or live transactions are involved. Proper attribution is
given to IBC standard authors, process algebra pioneers, and all cited
researchers.

\section{Summary}\label{summary}

This chapter has presented the conceptual framework and methodology for
analysing IBC\textquotesingle s asset and data behaviours. The key
elements are:

\begin{itemize}
    \item \textbf{Conceptual framework:} Operational definitions of observable behaviour, traces, branching structure, and deadlock specifically for IBC; three analytical scenarios (Unified, Stacked, Conflicting) with formal criteria; conceptual diagram showing the analysis flow
    \item \textbf{Methodology:} Five phases: (1) domain analysis and atomic action identification, (2) process algebra specification, (3) LTS extraction, (4) behavioural comparison using weak bisimulation and deadlock detection, (5) documentation for reusability
    \item \textbf{Validation:} Correctness through face validity and systematic mapping; completeness through state space coverage; reproducibility through complete specification in appendix
\end{itemize}
